\documentclass[11pt]{amsart}

\usepackage{macros}
\usepackage{refcount,xstring}
\usepackage{bbm}
\usepackage{subfiles}

%\makeatletter
%\numberwithin{section}{chapter}
%\def\@secnumfont{\mdseries}
%\def\section{\@startsection{section}{1}%
%  \z@{.7\linespacing\@plus\linespacing}{.5\linespacing}%
%  {\normalfont\scshape\centering}}
%\def\subsection{\@startsection{subsection}{2}%
%  \z@{.5\linespacing\@plus.7\linespacing}{-.5em}%
%  {\normalfont\bfseries}}

%\patchcmd{\@thm}{\let\thm@indent\indent}{\let\thm@indent\noindent}{}{}
%\patchcmd{\@thm}{\thm@headfont{\scshape}}{\thm@headfont{\bfseries}}{}{}

\newcommand\thfd{\d'}

\addbibresource{sections/refs}

\title{Poisson structures and derived spheres}

\begin{document}

\maketitle
\tableofcontents


\section{Non-topological Poisson sigma models}
\subfile{sections/aksz}	

\section{Manin pairs, Manin triples, and Poisson--Lie groups}

We first recall some standard definitions.

\begin{dfn}[{\cite[\S3]{Drinfeld}, \cite[Definition 2.1.1]{AK-S}, among others}]
    Let $\fd$ be a real or complex Lie algebra equipped with a nondegenerate symmetric invariant pairing, of split signature in the real case. \begin{itemize}
            \item A \emph{Manin pair} consists of $\fd$ together with a choice of maximal isotropic Lie subalgebra $\fg_+ \subset \fd$. 
            \item A \emph{Manin quasi-triple} consists of a Manin pair, together with a choice of complementary maximal isotropic linear subspace $V_- \subset \fd$.
            \item A \emph{Manin triple} consists of a Manin quasi-triple for which $V_- = \fg_-$ is a Lie subalgebra. %so that $\fd$ is the direct sum of 
            \item A \emph{Harish-Chandra Manin triple} consists of a Manin triple $\fg_+ \rightarrow \fd \leftarrow \fg_-$, together with a Lie group $G$ exponentiating $\fg_-$ and a $G$-representation on~$\fd$ extending the infinitesimal adjoint action of~$\fg_-$.
        \end{itemize}
\end{dfn}

\section{Cigar compactifications via homotopy transfer}
\subfile{sections/compactify}

\section{Three-sphere compactification}
\subfile{sections/m2}		


%\subfile{sections/pseudo}

%\subfile{sections/spheres}

\printbibliography

\end{document}
